% =====================================================================
%  CHAPTER 4: 特殊人群营养
% =====================================================================
\chapter{特殊人群营养}
\setcounter{chapter}{4}
\chapterintro{
\keyword{掌握}生命早期1000天机遇窗口及DOHaD理论；掌握孕妇、乳母、婴幼儿、学龄前儿童、学龄儿童、青少年、老年人的生理特点、营养需要及膳食指南关键推荐。\keyword{熟悉}运动员营养代谢特点及营养需要；熟悉各特殊人群常见营养问题及预防。\keyword{了解}高、低温环境人员营养需要和膳食原则。
}

% =====================================================================
\section{生命早期营养与DOHaD理论}
% =====================================================================

\begin{tipbox}
\textbf{DOHaD理论（Developmental Origins of Health and Disease，健康与疾病的发育起源）：}
生命早期（胎儿、婴幼儿期）经历不利因素（如营养不良、环境不良等），将会增加其成年后患肥胖、糖尿病、心血管疾病等慢性疾病的几率。其核心机制是"发育程序化（Developmental Programming）"——早期营养环境通过表观遗传调控（DNA甲基化、组蛋白修饰等）改变基因表达模式和组织分化方向，对个体远期疾病风险产生持久性影响。

\textbf{关键机制——"节俭表型"假说（Thrifty Phenotype Hypothesis）：}
宫内营养不良使胎儿产生"预测性适应"——调整代谢和内分泌系统以应对"匮乏"环境。若出生后实际面临"营养过剩"环境，则代谢-环境不匹配将大幅增加成年期代谢综合征风险。

\textbf{生命早期1000天（The First 1000 Days of Life）：}
生命早期1000天关键期 = 孕期270天 + 0-6月龄180天（母乳喂养）+ 7-24月龄（辅食添加）。这一阶段是生长发育的关键期和敏感期，此时营养状态对个体终身的体格发育、认知潜能、免疫功能和慢性疾病风险产生深远和不可逆的影响，是进行营养干预以获得最大健康收益的"机遇窗口期"。
\end{tipbox}

\begin{exambox}
\textbf{考点提示：} (1) 1000天的计算——孕期270d + 0-6月龄180d + 7-24月龄；注意不等于简单地"出生后730天"；(2) DOHaD的核心逻辑——早期环境 → 发育程序化（表观遗传）→ 远期健康/疾病；(3) "不匹配"是DOHaD解释代谢综合征的核心机制——宫内"预测匮乏" + 出生后"实际营养过剩" → 代谢-环境失匹配 → 代谢综合征高风险。
\end{exambox}

% =====================================================================
\section{孕妇营养}
% =====================================================================

\subsection{孕期生理特点}

\begin{keybox}
\textbf{妊娠期——生命早期1000天机遇窗口的起始阶段。}

\textbf{1. 内分泌变化：}
\begin{itemize}[leftmargin=*]
    \item \textbf{人绒毛膜促性腺激素（HCG）：}受精1周开始出现，8-9周达高峰。由胎盘合体滋养层细胞分泌，维持妊娠早期所需的孕酮分泌。
    \item \textbf{人绒毛膜生长激素（HCS）：}妊娠6周出现，36周达高峰至分娩。由胎盘合体滋养层细胞分泌，促进乳腺发育和胎儿生长，同时降低母体胰岛素敏感性。
    \item \textbf{雌激素：}分泌增加，增加子宫和胎盘的血流量，促进母体乳房发育。孕晚期可达非孕期的1000倍。
    \item \textbf{孕激素（孕酮）：}分泌增加，减少子宫收缩，维持妊娠。促进乳腺腺泡发育。
    \item \textbf{甲状腺激素：}水平升高 → 基础代谢率（BM）升高。
    \item \textbf{胰岛素敏感性下降：}由HCS、雌激素、孕激素等共同介导 → 妊娠期糖尿病（GDM）风险升高，属生理性胰岛素抵抗。
\end{itemize}

\textbf{2. 消化系统变化：}
\begin{itemize}[leftmargin=*]
    \item 贲门括约肌松弛 → 恶心呕吐（早孕反应）。
    \item 孕酮分泌增加 → 胃排空时间延长、胃肠蠕动减慢 → 胃肠胀气便秘。
    \item 对钙、铁、VitB$_{12}$、叶酸等营养素吸收增强——有利于营养储备。
\end{itemize}

\textbf{3. 循环系统变化：}
血容量增加，血液稀释：血浆容积增加幅度 $>$ 红细胞增加幅度 → \textbf{生理性贫血}。通常在孕20-30周最明显。

\textbf{WHO贫血标准（2024）：}
\begin{itemize}[leftmargin=*]
    \item 孕期：Hb $\geq$ 110 g/L
    \item 孕中期：Hb $\geq$ 105 g/L
\end{itemize}
\end{keybox}

\subsection{孕期体重增长}

\begin{table}[h!]
\centering
\caption{孕期推荐体重增长范围}
\begin{tabular}{lcc}
\hline
\textbf{孕前BMI分类} & \textbf{总增长范围（kg）} & \textbf{中晚期每周增长（kg/周）} \\
\hline
低体重（BMI $<$ 18.5） & 11.0-16.0 & 0.46（0.37-0.56） \\
正常体重（BMI 18.5-24.0） & 8.0-14.0 & 0.37（0.26-0.48） \\
超重（BMI 24.0-28.0） & 7.0-11.0 & 0.30（0.22-0.37） \\
肥胖（BMI $\geq$ 28.0） & 5.0-9.0 & 0.22（0.15-0.30） \\
\hline
\end{tabular}
\end{table}

\subsection{孕期营养需要}

\begin{keybox}
\textbf{孕期各阶段能量与宏量营养素增加量：}

\begin{center}
\begin{tabular}{lcccc}
\hline
\textbf{营养素} & \textbf{非孕期} & \textbf{孕早期} & \textbf{孕中期} & \textbf{孕晚期} \\
\hline
能量EER（kcal/d） & 1700 & +0 & +250 & +400 \\
蛋白质（g/d） & -- & +0 & +15 & +30 \\
\hline
\end{tabular}
\end{center}

\textbf{脂类：}供能比20\%-30\%不变。亚油酸占总能量4\%，$\alpha$-亚麻酸0.6\%，EPA+DHA达到250 mg/d。

\textbf{碳水化合物：}孕早期不低于130 g/d——预防酮体对胎儿中枢神经系统的毒性。

\textbf{关键微量营养素：}

\begin{center}
\begin{tabular}{lccc}
\hline
\textbf{营养素} & \textbf{非孕期} & \textbf{孕期RNI} & \textbf{备注} \\
\hline
钙（mg/d） & 800 & 800（不增加） & 孕晚期胎儿大量储存，钙吸收率代偿性升高 \\
铁（mg/d） & 18 & 孕中25，孕晚29 & 用于母体红细胞增加、胎儿胎盘需要、产时失血储备 \\
锌（mg/d） & 8.5 & 10.5 & -- \\
碘（$\mu$g/d） & 120 & 230 & 严重缺乏 → 克汀病 \\
叶酸（$\mu$g DFE/d） & -- & 400 & 备孕3个月起补充，持续整个妊娠期 \\
维生素D（$\mu$g/d） & 10 & 10（400 IU/d） & 缺乏 → 新生儿低钙血症和佝偻病风险增加 \\
\hline
\end{tabular}
\end{center}
\end{keybox}

\begin{keybox}
\textbf{叶酸——最重要的孕前/孕期微量营养素：}
\begin{itemize}[leftmargin=*]
    \item \textbf{推荐：}备孕妇女从计划怀孕前3个月开始每日补充叶酸\textbf{400 $\mu$g DFE/d}，持续整个妊娠期。
    \item \textbf{意义：}预防神经管缺陷（NTDs）——叶酸缺乏是胎儿神经管未能正常闭合（如脊柱裂、无脑儿）的确定病因。充足补充可使NTDs发生率降低70\%。
    \item \textbf{为什么必须从孕前开始？}孕后3-8周是神经管闭合关键期，此时许多妇女尚不知已怀孕。
    \item \textbf{膳食来源：}深绿色叶菜、肝脏、豆类、强化面粉。
\end{itemize}
\end{keybox}

\subsection{孕期膳食指南（2022版）}

\begin{keybox}
\textbf{《备孕和孕期妇女膳食指南（2022年）》6条关键推荐：}
\begin{enumerate}[leftmargin=*]
    \item \textbf{调整孕前体重至正常范围，保证孕期体重适宜增长。}
    \item \textbf{常吃含铁丰富的食物，选用碘盐，合理补充叶酸和维生素D。}
    \item \textbf{孕吐严重者，可少量多餐，保证摄入含必需量碳水化合物的食物。}
    \item \textbf{孕中晚期适量增加奶、鱼、禽、蛋、瘦肉的摄入。}
    \item \textbf{经常户外活动，禁烟酒，保持健康生活方式。}
    \item \textbf{愉快孕育新生命，积极准备母乳喂养。}
\end{enumerate}
\end{keybox}

\subsection{孕期常见营养问题及预防}

\begin{defbox}
\textbf{孕期常见营养问题：}

\begin{enumerate}[leftmargin=*]
    \item \textbf{早孕反应（妊娠剧吐 hyperemesis gravidarum）：}
    \begin{itemize}[leftmargin=*]
        \item 发生机制：HCG升高 + 贲门括约肌松弛 + 胃肠蠕动减慢。
        \item 营养风险：严重者进食困难 → 能量和营养素摄入不足 → 酮体产生（对胎儿中枢神经系统有毒性）。
        \item 预防处理：少量多餐（每日6-8次）；保证碳水化合物最低摄入量130 g/d（严重者需静脉补充葡萄糖，防止酮症）；晨起前吃干食物（如苏打饼干）；避免油腻和辛辣食物；补充维生素B$_{6}$有一定效果。
    \end{itemize}

    \item \textbf{孕期便秘：}
    \begin{itemize}[leftmargin=*]
        \item 发生机制：孕酮分泌增加 → 胃肠蠕动减慢；增大的子宫压迫直肠。
        \item 预防：增加膳食纤维（新鲜蔬果、全谷物）和足够饮水，适量运动。
    \end{itemize}

    \item \textbf{妊娠期糖尿病（GDM）：}
    \begin{itemize}[leftmargin=*]
        \item 发生机制：HCS、雌激素、孕激素等综合作用 → 生理性胰岛素抵抗。
        \item 发生率约10\%-20\%。高危因素：孕前超重/肥胖、有糖尿病家族史、高龄妊娠。
        \item 危害：增加巨大儿、剖宫产率、新生儿低血糖风险，母亲远期2型糖尿病风险升高。
        \item 预防：合理控制孕期体重增长，避免高GI食物和含糖饮料，适量运动。
    \end{itemize}

    \item \textbf{妊娠期高血压疾病：}
    \begin{itemize}[leftmargin=*]
        \item 与钙摄入不足可能有关，保证充足钙摄入（800 mg/d）可降低高血压风险。
        \item 控制钠盐摄入，维持适宜体重增长。
    \end{itemize}

    \item \textbf{孕期贫血（缺铁性贫血最常见）：}
    \begin{itemize}[leftmargin=*]
        \item 原因：血容量扩张 + 胎儿和胎盘铁转运 + 产时失血储备 → 铁需要量大幅增加。
        \item 孕中期铁RNI 25 mg/d，孕晚期29 mg/d（比非孕期18 mg/d显著增加）。
        \item 预防：常吃含铁丰富食物（动物肝脏、动物血、红肉），同时摄入维生素C促铁吸收。
    \end{itemize}

    \item \textbf{孕期钙营养与骨骼健康：}
    \begin{itemize}[leftmargin=*]
        \item 孕期钙RNI 800 mg/d（与孕前相同），因钙吸收率代偿性升高。
        \item 钙摄入严重不足时，胎儿动员母体骨骼钙 → 母体骨量损失。充足的奶及奶制品摄入可预防。
    \end{itemize}
\end{enumerate}
\end{defbox}

% =====================================================================

\subsection{泌乳量与母乳分期}

\begin{keybox}
\textbf{哺乳期每日泌乳量：}700-800 ml（平均750 ml）。

\textbf{母乳分期：}
\begin{enumerate}[leftmargin=*]
    \item \textbf{初乳（产后第1周）：}淡黄色，蛋白质多，富含免疫球蛋白（分泌型IgA和乳铁蛋白），乳糖和脂肪相对较少。
    \item \textbf{过渡乳（产后第2周）：}蛋白质$\downarrow$，乳糖和脂肪$\uparrow$。
    \item \textbf{成熟乳（第2周以后）：}乳白色，蛋白质相对较少，乳糖和脂肪多。
\end{enumerate}
\end{keybox}

\subsection{母乳喂养的优点}

\begin{keybox}
\textbf{母乳喂养的优点：}
\begin{enumerate}[leftmargin=*]
    \item \textbf{蛋白质：}乳清蛋白为主，必需氨基酸比例适当。
    \item \textbf{脂肪：}脂肪颗粒小 + 乳脂酶 → 易消化吸收。
    \item \textbf{糖类：}富含乳糖 → 促进乳酸杆菌、抑制大肠埃希菌生长，利于钙、铁、锌吸收。
    \item \textbf{矿物质：}钙磷比例适宜（2:1）。
    \item \textbf{含大量免疫物质：}分泌型IgA、乳铁蛋白、溶菌酶等。
    \item \textbf{不易过敏。}
    \item \textbf{经济、方便、卫生。}
    \item \textbf{促进产后恢复，增进母婴交流。}
\end{enumerate}
\end{keybox}

\subsection{乳母的营养需要}

\begin{keybox}
\textbf{乳母关键营养素较非孕妇女的增加量：}

\begin{center}
\begin{tabular}{lcl}
\hline
\textbf{营养素} & \textbf{增加量} & \textbf{备注} \\
\hline
能量EER & +400 kcal/d & 满足泌乳能量消耗 \\
蛋白质 & +25 g/d & 高品质蛋白质 \\
钙 & 800 mg/d（不增） & 乳汁中含量较稳定，每日排出300 mg \\
  & & 乳母钙供给不足会用自身骨骼中的钙满足乳汁钙含量 \\
铁 & -- & 铁不能通过乳腺输送到乳汁 \\
\hline
\end{tabular}
\end{center}

\textbf{重要说明：}
\begin{itemize}[leftmargin=*]
    \item \textbf{钙：}乳汁中钙含量较稳定（约300 mg/L），每日排出约300 mg。乳母钙供给不足时会动用自身骨骼中的钙来满足乳汁钙含量——哺乳期骨骼密度下降，断奶后可逆性恢复。
    \item \textbf{铁：}铁不能通过乳腺输送到乳汁，因此母乳中铁含量不受母体铁状态影响。
\end{itemize}
\end{keybox}

\subsection{哺乳期膳食指南}

\begin{keybox}
\textbf{哺乳期膳食指南关键推荐：}
\begin{enumerate}[leftmargin=*]
    \item 增加富含优质蛋白质及维生素A的动物性食物和海产品摄入。
    \item 保证充足的能量和液体摄入——能量在孕前基础上增加400 kcal/d。
    \item 增加钙摄入，重视哺乳期骨骼健康——每日饮奶400-500 ml，保证钙摄入800 mg/d。
    \item 保持身心愉快，保证充足睡眠——充足休息是维持泌乳量的必要条件。
    \item 坚持哺乳，适量身体活动——有助于产后体重恢复。
    \item 禁烟酒，避免浓茶和咖啡。
\end{enumerate}
\end{keybox}

\subsection{乳母常见营养问题及预防}

\begin{defbox}
\textbf{乳母常见营养问题：}

\begin{enumerate}[leftmargin=*]
    \item \textbf{能量和营养素摄入不足：}
    \begin{itemize}[leftmargin=*]
        \item 原因：产后急于恢复体型而过度控制饮食；照顾婴儿导致饮食不规律。
        \item 危害：泌乳量减少，乳汁质量降低，母体营养耗竭。
        \item 预防：能量在孕前基础上增加400 kcal/d，蛋白质增加25 g/d；保证充足液体摄入（多喝汤水，泌乳需水约1 L/d额外补充）。
    \end{itemize}

    \item \textbf{乳母骨骼健康问题：}
    \begin{itemize}[leftmargin=*]
        \item 乳汁每日排出约300 mg钙，钙供给不足时动用人骨钙 → 哺乳期骨密度暂时性下降（约3\%-5\%）。
        \item 预防：保证每日钙摄入800 mg/d（饮奶400-500 ml）；断奶后骨密度一般可恢复。
    \end{itemize}

    \item \textbf{铁营养与产后贫血：}
    \begin{itemize}[leftmargin=*]
        \item 铁不能通过乳腺进入乳汁，影响的是母亲自身的铁储备。
        \item 分娩失血 + 哺乳消耗可导致产后缺铁性贫血。预防：注意补铁（动物肝脏、红肉），产后继续补充铁剂至血象恢复。
    \end{itemize}

    \item \textbf{乳腺炎：}
    \begin{itemize}[leftmargin=*]
        \item 充分排空乳汁是最重要的预防手段之一。营养方面注意保证充足的蛋白质和维生素摄入（增强免疫力）。
    \end{itemize}

    \item \textbf{维生素缺乏：}
    \begin{itemize}[leftmargin=*]
        \item 维生素A：乳汁维生素A含量受乳母膳食维生素A影响较大，维生素A摄入不足导致乳汁含量降低。预防：摄入富含维生素A的动物性食物和海产品。
        \item 维生素D：乳母维生素D缺乏 → 乳汁维生素D含量低 → 婴儿维生素D不足。乳母需保证维生素D摄入（10 $\mu$g/d）。
        \item 维生素B$_{1}$和B$_{2}$需要量与能量摄入相关（0.5 mg/1000 kcal），需相应增加摄入。
    \end{itemize}
\end{enumerate}
\end{defbox}

% =====================================================================

\subsection{婴幼儿生理特点}

\begin{keybox}
\textbf{婴幼儿（0-3岁）主要生理特点：}

\begin{enumerate}[leftmargin=*]
    \item \textbf{生长发育旺盛：}第一生长高峰期——婴儿期体重增长为出生时的3倍（1岁时约10 kg），身高增加约25 cm。2岁至青春期前：体重每年增长约2 kg，身高5-7 cm。
    \item \textbf{代谢率高：}基础代谢率约为成人的2倍，能量需要量显著高于成人（按每千克体重计）。
    \item \textbf{消化系统不成熟：}胃容量小（新生儿约30-50 ml，1岁时250-350 ml），肠道通透性高，消化酶分泌不足——对食物的消化吸收能力有限，易发生消化紊乱。
    \item \textbf{神经系统：}大脑快速发育，3岁时脑重约为成人80\%，6岁时可达成人90\%。突触形成和髓鞘化在2岁内最为活跃，DHA和ARA充足供给对脑发育至关重要。
    \item \textbf{免疫功能：}新生儿免疫系统尚未发育完全，IgG可从母体获得（胎盘转运），sIgA主要来自母乳（初乳中含量极高）——母乳喂养对婴儿免疫保护至关重要。
    \item \textbf{肾功能：}肾小球滤过率和肾小管重吸收功能都不成熟，肾溶质负荷有限——不能处理高浓度废物（需大量水分），不宜摄入过多蛋白质和矿物质（尤其钠）。
    \item \textbf{体成分特点：}体内水分含量高（新生儿约75\%，成人约60\%），体表面积相对较大 → 水分蒸发量大 → 每日需水量150 ml/(kg$\cdot$d)。
    \item \textbf{铁贮存：}足月新生儿体内贮存约300 mg铁，可供生后4-6个月需要。6月龄后贮存铁耗竭，必须从辅食补充。
\end{enumerate}
\end{keybox}

\subsection{0-6月龄婴儿——纯母乳喂养}

\begin{keybox}
\textbf{《0-6月龄婴儿母乳喂养指南（2022）》准则：}
\begin{enumerate}[leftmargin=*]
    \item \textbf{准则1：母乳是婴儿最理想的食物，坚持6月龄内纯母乳喂养。}
    \item \textbf{准则2：生后1小时内开奶。}
    \item \textbf{准则3：回应式喂养。}
    \item \textbf{准则4：适当补充维生素D，母乳喂养无需补钙。}
    \item \textbf{准则5：任何动摇母乳喂养的想法和举动都必须咨询医生。}
    \item \textbf{准则6：定期监测体格指标。}
\end{enumerate}
\end{keybox}

\subsection{7-24月龄婴幼儿——辅食添加}

\begin{keybox}
\textbf{《7-24月龄婴幼儿喂养指南（2022）》准则：}
\begin{enumerate}[leftmargin=*]
    \item \textbf{准则1：继续母乳喂养，满6月龄起必须添加辅食，从富含铁的泥糊状食物开始。}
    \item \textbf{准则2：及时引入多样化食物，重视动物性食物的添加。}
    \item \textbf{准则3：尽量少加糖、盐，油脂适当，保持食物原味。}
    \item \textbf{准则4：提倡回应式喂养，鼓励但不强迫进食。}
    \item \textbf{准则5：注重饮食卫生和进食安全。}
    \item \textbf{准则6：定期监测体格指标，追求健康生长。}
\end{enumerate}
\end{keybox}

\begin{keybox}
\textbf{辅食添加的意义：}
\begin{enumerate}[leftmargin=*]
    \item \textbf{补充营养素：}6月龄后母乳已不能满足婴儿能量和营养素（尤其铁、锌）的全部需要——\textbf{7-12月龄99\%的铁必须从辅食中获得}。
    \item \textbf{促进消化系统发育：}辅食刺激胃肠道消化酶分泌和功能成熟。
    \item \textbf{培养咀嚼和吞咽能力：}不同质地食物的逐步引入有利于口腔运动和咀嚼功能发展。
    \item \textbf{促进味觉发育和饮食习惯形成：}辅食添加期是接受新食物和建立食物偏好的关键窗口期。
    \item \textbf{促进神经心理发育：}自主进食习惯的培养有利于精细运动和认知发育。
\end{enumerate}
\end{keybox}

\begin{keybox}
\textbf{辅食添加的原则：}
\begin{itemize}[leftmargin=*]
    \item 每次只添加一种新食物，观察3-5天以排除过敏。
    \item 由少到多，由稀到稠，由细到粗。
    \item 从泥糊状 → 碎末状 → 小块状 → 条状。
    \item 第一口辅食首选\textbf{强化铁的米粉、动物肝脏泥、瘦肉泥}等富含铁的泥糊状食物。
    \item 应在婴儿健康、消化功能正常时添加。
    \item 保持原味，不加盐（1岁前）、糖及刺激性调味品；1岁后逐渐尝试淡口味家庭膳食。
    \item 6月龄时婴儿体内储存的铁已耗尽——\textbf{7-12月龄99\%的铁必须从辅食中获得}。
\end{itemize}
\end{keybox}

\subsection{婴幼儿常见营养问题及预防}

\begin{defbox}
\textbf{婴幼儿常见营养问题：}

\begin{enumerate}[leftmargin=*]
    \item \textbf{缺铁性贫血（IDA）：}
    \begin{itemize}[leftmargin=*]
        \item 6月龄至2岁为高发期——原因：6月后贮存铁耗竭，辅食铁来源不足。
        \item 7-12月龄铁RNI = 10 mg/d，必须通过辅食获得。
        \item 危害：影响认知发育和行为能力（即使纠正后，认知缺陷可能持续）。
        \item 预防：保证6月龄及时添加富含铁辅食（强化铁米粉、肝脏泥、肉泥）；适当补充维生素C促进铁吸收。
    \end{itemize}

    \item \textbf{维生素D缺乏性佝偻病：}
    \begin{itemize}[leftmargin=*]
        \item 原因：母乳中维生素D含量极低，日照不足，未补充维生素D制剂。
        \item 预防：生后数日开始每日补充维生素D 400 IU（10 $\mu$g/d），至2岁。
    \end{itemize}

    \item \textbf{蛋白质-能量营养不良（PEM）：}
    \begin{itemize}[leftmargin=*]
        \item 辅食添加不及时、辅食质量差、断奶后喂养不当。
        \item 表现：生长迟缓（身高别年龄低于标准）、消瘦（体重别身高低于标准）。
        \item 预防：保证辅食的质和量，及时引入动物性食物。
    \end{itemize}

    \item \textbf{锌缺乏：}
    \begin{itemize}[leftmargin=*]
        \item 表现：生长发育迟缓、食欲不振、异食癖、免疫功能下降。
        \item 预防：辅食中添加富含锌的食物（红肉、肝脏、海产品）。
    \end{itemize}

    \item \textbf{碘缺乏：}
    \begin{itemize}[leftmargin=*]
        \item 严重缺乏可致克汀病（先天性碘缺乏综合征：矮小、智力严重障碍）。
        \item 预防：使用碘盐。
    \end{itemize}

    \item \textbf{维生素K缺乏：}
    \begin{itemize}[leftmargin=*]
        \item 母乳中维生素K含量很低，单纯母乳喂养新生儿易缺乏 → 出血性疾病风险。
        \item 预防：新生儿出生后立即注射维生素K（肌注1 mg）。
    \end{itemize}

    \item \textbf{龋齿：}
    \begin{itemize}[leftmargin=*]
        \item 含糖饮料（含果汁、加糖乳品）和甜食是婴幼儿龋齿的主要膳食因素。
        \item 预防：1岁前不加糖，不喂含糖饮料，1岁后减少游离糖摄入。
    \end{itemize}
\end{enumerate}
\end{defbox}

\subsection{婴幼儿营养需要}

\begin{keybox}
\textbf{婴幼儿关键营养素需要特点：}

\begin{itemize}[leftmargin=*]
    \item \textbf{蛋白质：}优质蛋白占总蛋白质1/2。
    \item \textbf{钙/磷比值：}适宜比值（2.3:1），比牛奶（1.4:1）合理。
    \item \textbf{铁：}足月新生儿有足够铁贮存，供4-6个月需要。\textbf{6m-2y易患缺铁性贫血}。7-12月龄RNI = 10 mg/d。
    \item \textbf{锌：}婴幼儿期缺锌 → 生长发育迟缓、脑发育受损、食欲不振、异食癖。
    \item \textbf{碘：}缺乏 → 克汀病（先天性碘缺乏综合征：矮小、智力严重障碍）。
    \item \textbf{维生素D：}乳类含量微量，需要额外补充（生后数日开始补充400 IU/d）。
    \item \textbf{维生素K：}母乳中很少，单纯母乳喂养易缺乏。
\end{itemize}
\end{keybox}

% =====================================================================
\section{学龄前儿童营养}
% =====================================================================

\subsection{生理特点}

\begin{defbox}
\textbf{学龄前儿童（3-6岁）主要生理特点：}

\begin{enumerate}[leftmargin=*]
    \item \textbf{生长发育：}每年体重增长约2 kg，每年身高增长约5-7 cm。生长速率较婴幼儿期明显趋缓。
    \item \textbf{神经系统：}依旧属于脑细胞生长发育关键时期。3岁脑重约为成人80\%，6岁可达成人90\%。
    \item \textbf{咀嚼和消化能力仍有限：}3岁乳牙已出齐，6岁恒牙已萌出，但咀嚼和消化能力远低于成人。
    \item \textbf{心理发育特点：}注意力分散（约15分钟），无法专心进食。好奇心强，模仿能力增强，此时期是培养良好饮食行为习惯的关键窗口期。
    \item \textbf{常见营养问题：}挑食偏食、零食摄入过多、含糖饮料摄入过量、钙和铁摄入不足。
\end{enumerate}
\end{defbox}

\subsection{学龄前儿童膳食指南（2022版）}

\begin{keybox}
\textbf{学龄前儿童膳食指南（2022）核心推荐：}
\begin{enumerate}[leftmargin=*]
    \item \textbf{食物多样，规律就餐，自主进食}——培养良好饮食习惯。
    \item \textbf{每天饮奶，足量饮水，合理选择零食。}
    \begin{itemize}[leftmargin=*]
        \item 奶类每日300-500 mL。
        \item 水每日600-800 mL（以白开水为主）。
    \end{itemize}
    \item \textbf{合理烹调，少调料少油炸}——采用蒸、煮、炖、煨等方式。每人每日食盐 $<$ 3 g。
    \item \textbf{参与食物选择制作}——增进对食物的认知与喜爱。
    \item \textbf{经常户外活动，定期体格测量。}
    \begin{itemize}[leftmargin=*]
        \item 每天至少60分钟体育活动。
        \item 每天看电视、玩平板电脑累计时间不超过2小时。
    \end{itemize}
\end{enumerate}
\end{keybox}

\subsection{学龄前儿童常见营养问题及预防}

\begin{defbox}
\textbf{学龄前儿童常见营养问题：}

\begin{enumerate}[leftmargin=*]
    \item \textbf{挑食偏食：}
    \begin{itemize}[leftmargin=*]
        \item 学龄前儿童自主意识增强，对陌生食物有"恐新症"。
        \item 预防：家长以身作则，树立良好榜样；多次尝试（8-15次）让儿童接受新食物；不强迫进食，营造愉快用餐氛围；让儿童参与食物选择和制作。
    \end{itemize}

    \item \textbf{零食摄入过多：}
    \begin{itemize}[leftmargin=*]
        \item 劣质零食（高能量、高脂肪、高盐、高糖、低营养）挤占主食和奶类摄入。
        \item 预防：合理选择零食（奶制品、水果、坚果），控制零食量和时间（两餐之间，不代替正餐）。
    \end{itemize}

    \item \textbf{含糖饮料摄入过量：}
    \begin{itemize}[leftmargin=*]
        \item 含糖饮料 → 能量过剩 → 超重/肥胖 + 龋齿风险。
        \item 预防：以白开水为主要饮品（每日600-800 ml），不提供含糖饮料。
    \end{itemize}

    \item \textbf{钙摄入不足：}
    \begin{itemize}[leftmargin=*]
        \item 奶类摄入不足是最常见原因 → 影响骨骼发育和峰值骨量积累。
        \item 预防：保证每日奶类300-500 ml。
    \end{itemize}

    \item \textbf{铁和锌摄入不足：}
    \begin{itemize}[leftmargin=*]
        \item 动物性食物摄入不足的儿童易缺铁和锌。
        \item 预防：保证肉类、肝脏、海产品适量摄入。
    \end{itemize}

    \item \textbf{超重和肥胖：}
    \begin{itemize}[leftmargin=*]
        \item 能量摄入大于消耗，零食和含糖饮料过多，体育活动不足。
        \item 预防：控制总能量和零食，保证每日户外活动至少60分钟。
    \end{itemize}
\end{enumerate}
\end{defbox}

% =====================================================================
\section{学龄儿童营养}
% =====================================================================

\subsection{生理特点}

\begin{defbox}
\textbf{学龄儿童（6-12岁）主要生理特点：}

\begin{enumerate}[leftmargin=*]
    \item \textbf{生长发育：}身高和体重较快增长但低于学龄前期和青春期。
    \item \textbf{骨骼：}骨胶多、钙盐少弹性大、硬度小，不易骨折而易弯曲变形。
    \item \textbf{换牙：}恒牙6-7岁；钙质沉积更早。
    \item \textbf{神经系统：}脑细胞分化基本完成，神经传导逐渐迅速、准确。
    \item \textbf{心理特点：}控制力仍较弱；精神易紧张，情绪易波动，过多的脑力活动易疲劳。
\end{enumerate}

\textbf{生长发育4规律：}
\begin{enumerate}[leftmargin=*]
    \item 连续性
    \item 不同步
    \item 波浪式
    \item 差异性
\end{enumerate}
\end{defbox}

\subsection{学龄儿童膳食指南（2022版）}

\begin{keybox}
\textbf{学龄儿童膳食指南（2022）5条核心推荐：}
\begin{enumerate}[leftmargin=*]
    \item \textbf{主动参与食物制作和选择，提高营养素养。}
    \item \textbf{吃好早餐，合理选择零食，培养健康饮食行为。}
    \item \textbf{天天喝奶，足量饮水，不喝含糖饮料，禁止饮酒。}
    \item \textbf{多户外活动，少视屏时间，每天60分钟以上的中高强度身体活动。}
    \item \textbf{定期监测体格发育，保持体重适宜增长。}
\end{enumerate}
\end{keybox}

\subsection{学龄儿童常见营养问题及预防}

\begin{defbox}
\textbf{学龄儿童常见营养问题：}

\begin{enumerate}[leftmargin=*]
    \item \textbf{不吃早餐或早餐质量差：}
    \begin{itemize}[leftmargin=*]
        \item 学龄儿童不吃早餐比例较高，影响上午学习效率和认知功能。
        \item 预防：培养规律早餐习惯；早餐应包含谷类、奶类/蛋类/肉类和蔬菜水果三类食物以上。
    \end{itemize}

    \item \textbf{零食和含糖饮料过度摄入：}
    \begin{itemize}[leftmargin=*]
        \item 校门口摊贩和自动售货机的零食饮料影响正餐摄入。
        \item 预防：合理选择零食，不喝含糖饮料，以白开水为日常饮水。
    \end{itemize}

    \item \textbf{钙和维生素D不足：}
    \begin{itemize}[leftmargin=*]
        \item 奶类摄入不足 + 日照减少（室内活动多） → 钙和维生素D缺乏风险。
        \item 预防：天天喝奶（300-500 ml/d），增加户外活动。
    \end{itemize}

    \item \textbf{超重肥胖或消瘦两极分化：}
    \begin{itemize}[leftmargin=*]
        \item 超重肥胖：高能量密度食物过多，身体活动不足。
        \item 消瘦：偏食、厌食、对体型的不正确认知。
        \item 预防：平衡膳食 + 每天至少60分钟中高强度身体活动 + 定期监测体格发育。
    \end{itemize}

    \item \textbf{脊柱和骨骼发育问题：}
    \begin{itemize}[leftmargin=*]
        \item 学龄儿童骨骼骨胶多、钙盐少、硬度小，易弯曲变形。
        \item 姿势不良（长时间不正确的读写姿势） → 脊柱侧弯、驼背。
        \item 预防：保证钙摄入，保持正确坐姿，适量运动。
    \end{itemize}
\end{enumerate}
\end{defbox}

% =====================================================================
\section{青少年营养}
% =====================================================================

\subsection{生理特点}

\begin{defbox}
\textbf{青少年——第二生长发育高峰期。}

\textbf{1. 体格发育：}
\begin{itemize}[leftmargin=*]
    \item \textbf{身高性别差异：}男较女晚2年；男7-9 cm/y；女5-7 cm/y。
    \item \textbf{体重：}4-5 kg/y，时间较长。
    \item \textbf{形态：}四肢比脊柱早；脚最先；从远到近。
\end{itemize}

\textbf{2. 性发育：}10岁前处于静止状态，青春期后在激素作用下迅速发育。

\textbf{3. 骨量积累：}青春期是获得峰值骨量（Peak Bone Mass）的最关键阶段——约45\%的成人骨量在青春期积累。

\textbf{4. 认知和心理：}抽象思维能力增强，独立意识增加，同伴影响力显著上升。易出现饮食行为偏离（节食减肥、不吃早餐、在外就餐增多）。
\end{defbox}

\subsection{青少年营养需要}

\begin{keybox}
\textbf{青少年关键营养素需求特点：}
\begin{itemize}[leftmargin=*]
    \item \textbf{能量：}显著增加。11-13岁：男2050 kcal/d，女1800 kcal/d；14-17岁：男2500 kcal/d，女2000 kcal/d。
    \item \textbf{蛋白质：}55-75 g/d（男 $>$ 女），优质蛋白质占1/2以上。长期摄入不足 → 青春期生长迟滞。
    \item \textbf{钙：}1000-1200 mg/d——峰值骨量积累关键窗口期。摄入不足 → 无法达到遗传决定的峰值骨量 → 老年期骨质疏松风险大幅增加。
    \item \textbf{铁：}青春期需求量大幅增加。男孩来自肌肉和血容量增加，女孩需弥补月经铁丢失。11-13岁：男15 mg/d，女18 mg/d；14-17岁：男16 mg/d，女18 mg/d。\textbf{青春期女性是缺铁性贫血高发人群。}
    \item \textbf{锌：}9-11.5 mg/d。缺乏 → 青春期生长迟缓和性发育延迟。
    \item \textbf{碘：}110-120 $\mu$g/d。缺乏 → 甲状腺肿和认知功能下降。
    \item \textbf{维生素D：}10 $\mu$g/d（400 IU/d）。
\end{itemize}
\end{keybox}

\subsection{青少年常见营养问题及预防}

\begin{defbox}
\textbf{青少年常见营养问题：}

\begin{enumerate}[leftmargin=*]
    \item \textbf{青春期女性缺铁性贫血：}
    \begin{itemize}[leftmargin=*]
        \item 原因：月经铁丢失 + 生长发育血容量扩大双重要求 → 是缺铁性贫血高发人群。
        \item 青少年女性铁RNI = 18 mg/d。预防：常吃含铁丰富食物（动物肝脏、动物血、红肉），同时摄入维生素C促铁吸收。
    \end{itemize}

    \item \textbf{钙摄入不足与峰值骨量受影响：}
    \begin{itemize}[leftmargin=*]
        \item 青春期是获得峰值骨量的最关键阶段——约45\%的成人骨量在此期积累。
        \item 钙摄入不足 → 峰值骨量无法达到遗传决定的最大值 → 老年期骨质疏松风险大幅增加。
        \item 预防：奶及奶制品为钙最佳来源，钙RNI 1000-1200 mg/d。
    \end{itemize}

    \item \textbf{不健康减肥行为：}
    \begin{itemize}[leftmargin=*]
        \item 盲目节食、不吃主食、过度运动、使用减肥药等，女性多见。
        \item 后果：能量和营养素不足，影响生长发育，月经紊乱/闭经。
    \end{itemize}

    \item \textbf{不吃早餐和在外就餐增多：}
    \begin{itemize}[leftmargin=*]
        \item 学习压力、时间紧张 → 不吃早餐率高。
        \item 在外就餐增多 → 高油高盐高能量膳食 + 蔬果摄入不足 → 超重肥胖风险增加。
    \end{itemize}

    \item \textbf{超重肥胖：}
    \begin{itemize}[leftmargin=*]
        \item 能量摄入过剩 + 缺乏运动 + 含糖饮料摄入过多。
        \item 预防：控制总能量，增加体力活动（每天至少60分钟中高强度运动），减少屏幕时间。
    \end{itemize}

    \item \textbf{锌和碘缺乏：}
    \begin{itemize}[leftmargin=*]
        \item 锌缺乏 → 青春期生长迟缓和性发育延迟（锌RNI 9-11.5 mg/d）。
        \item 碘缺乏 → 甲状腺肿和认知功能下降（碘RNI 110-120 $\mu$g/d）。
    \end{itemize}

    \item \textbf{情绪化进食与饮食行为偏离：}
    \begin{itemize}[leftmargin=*]
        \item 青春期内分泌波动 + 心理压力 → 情绪化进食（暴饮暴食）、厌食。
        \item 同伴影响力上升 → 饮食行为易受同伴影响（偏食、节食模仿等）。
    \end{itemize}
\end{enumerate}
\end{defbox}

% =====================================================================
\section{老年人营养}
% =====================================================================

\subsection{老年人生理特点与少肌症}

\begin{defbox}
\textbf{老年人主要生理变化：}

\textbf{1. 体成分变化——少肌症（Sarcopenia）：}
\begin{itemize}[leftmargin=*]
    \item 50岁以后，骨骼肌量平均每年减少1\%-2\%。
    \item 60岁以上慢性肌肉丢失估计约30\%，80岁以上为50\%。
    \item 表现为骨骼肌质量、力量和功能的进行性丧失。后果：跌倒、骨折、失能、生活质量下降和死亡率增加。
    \item 体内水分和瘦组织减少，体脂含量增加（尤其向心性脂肪分布）。
\end{itemize}

\textbf{2. 代谢变化：}
\begin{itemize}[leftmargin=*]
    \item 基础代谢率（BMR）下降：约每10年下降2\%。
    \item 分解代谢 $>$ 合成代谢：常出现负氮平衡。
\end{itemize}

\textbf{3. 消化系统功能减弱：}
\begin{itemize}[leftmargin=*]
    \item 胃酸分泌减少（萎缩性胃炎高发 → 影响VitB$_{12}$和钙吸收）。
    \item 肠蠕动减慢 → 便秘。
    \item 味觉（味蕾数量减少）和嗅觉敏感度下降 → 食欲降低。
\end{itemize}
\end{defbox}

\subsection{老年人营养需要}

\begin{keybox}
\textbf{老年人关键营养素RNI：}

\begin{center}
\begin{tabular}{lcc}
\hline
\textbf{营养素} & \textbf{50-64岁} & \textbf{65岁以上} \\
\hline
蛋白质（男/女，g/d） & 65 / 55 & 72 / 62 \\
钙（mg/d） & 800 & 800 \\
维生素D（$\mu$g/d） & 10 & 15 \\
\hline
\end{tabular}
\end{center}

\textbf{重要说明：}
\begin{enumerate}[leftmargin=*]
    \item \textbf{蛋白质需要量不减反增：}老年人蛋白质推荐摄入量为\textbf{1.0-1.5 g/(kg$\cdot$bw)/d}以预防少肌症。65岁以上：男72 g/d，女62 g/d——均高于50岁组的65 g/d和55 g/d。原因：合成抵抗（Anabolic Resistance）——衰老肌肉对较低剂量氨基酸/蛋白质的合成代谢反应下降，需要更多蛋白质刺激同等幅度的肌蛋白合成。
    \item \textbf{能量需求降低：}BMR下降和身体活动减少。
    \item \textbf{脂肪：}供能比20\%-30\%。
    \item \textbf{碳水化合物：}供能比50\%-65\%，选择低GI食物。
    \item \textbf{维生素D：}65岁以上RNI = 15 $\mu$g/d（成人10 $\mu$g/d）。因皮肤合成能力下降和日照减少。
    \item \textbf{水：}老年人渴觉敏感性下降，易发生隐性脱水。每日饮水量不少于1500 ml。
\end{enumerate}
\end{keybox}

\subsection{老年人膳食指南（2022版）}

\begin{keybox}
\textbf{《中国老年人膳食指南》条目（2022）：}
\begin{enumerate}[leftmargin=*]
    \item \textbf{食物品种丰富，动物性食物充足，常吃大豆制品。}
    \item \textbf{鼓励共同进餐，保持良好食欲。}
    \item \textbf{积极户外活动，延缓肌肉衰减，保持适宜体重。}
    \begin{itemize}[leftmargin=*]
        \item 结合抗阻运动（举轻哑铃、弹力带训练）最大程度对抗少肌症。
        \item BMI应不低于20 kg/m$^{2}$，不宜过度追求"瘦"。
    \end{itemize}
    \item \textbf{定期健康体检，测评营养状况，预防营养缺乏。}
\end{enumerate}
\end{keybox}

\begin{defbox}
\textbf{成功衰老（Successful Aging）：}在老年期尽可能延长健康寿命，维持高水平的生理和认知功能，仅于生命末期极短时间段出现功能下降和失能。三大要素（Rowe \& Kahn模型）：(1) 避免疾病和疾病相关失能；(2) 维持高水平的认知和身体功能；(3) 积极投入生活。
\end{defbox}

\subsection{老年人常见营养问题及预防}

\begin{defbox}
\textbf{老年人常见营养问题：}

\begin{enumerate}[leftmargin=*]
    \item \textbf{少肌症（Sarcopenia）：}
    \begin{itemize}[leftmargin=*]
        \item 定义：随年龄增长的骨骼肌质量、力量和功能的进行性丧失。
        \item 50岁后骨骼肌量年均减少1\%-2\%，60岁以上慢性肌肉丢失约30\%，80岁以上约50\%。
        \item 后果：跌倒、骨折、失能、生活质量下降和死亡率增加。
        \item 预防：充足优质蛋白质（1.0-1.5 g/(kg$\cdot$d)）+ 抗阻运动（举轻哑铃、弹力带训练）。
    \end{itemize}

    \item \textbf{蛋白质-能量营养不良（PEM）：}
    \begin{itemize}[leftmargin=*]
        \item 原因：食欲降低（味觉嗅觉减退）、咀嚼困难（牙齿脱落）、购买和烹饪能力下降、独居进食单调。
        \item 预防：保证动物性食物充足，常吃大豆制品，鼓励共同进餐。
    \end{itemize}

    \item \textbf{骨质疏松和骨折：}
    \begin{itemize}[leftmargin=*]
        \item 骨量随年龄下降，雌激素减少（女性绝经后骨丢失加速）。
        \item 预防：充足钙（800-1000 mg/d）和维生素D（65岁以上15 $\mu$g/d），规律负重运动。
    \end{itemize}

    \item \textbf{维生素D缺乏：}
    \begin{itemize}[leftmargin=*]
        \item 皮肤合成能力下降 + 日照减少 → 维生素D缺乏率高。
        \item 预防：多晒太阳，必要时补充制剂（65岁以上RNI 15 $\mu$g/d）。
    \end{itemize}

    \item \textbf{维生素B$_{12}$缺乏：}
    \begin{itemize}[leftmargin=*]
        \item 萎缩性胃炎高发 → 胃酸分泌减少 → 食物结合态B$_{12}$释放和吸收障碍。
        \item 预防：保证动物性食物摄入（肉类、鱼类、蛋类、奶类），严重者需补充制剂。
    \end{itemize}

    \item \textbf{便秘：}
    \begin{itemize}[leftmargin=*]
        \item 肠蠕动减慢 + 膳食纤维摄入不足 + 饮水不足 + 活动减少。
        \item 预防：增加膳食纤维和饮水（不少于1500 ml/d），适量运动。
    \end{itemize}

    \item \textbf{隐性脱水：}
    \begin{itemize}[leftmargin=*]
        \item 原因：渴觉敏感性下降 → 不自觉饮水不足。
        \item 预防：提醒喝水，每日饮水量不少于1500 ml（心肾功能正常者）。
    \end{itemize}

    \item \textbf{牙齿和口腔问题：}
    \begin{itemize}[leftmargin=*]
        \item 牙齿脱落、咀嚼困难 → 固体食物摄入受限 → 蔬菜、肉类摄入不足 → 多种营养素缺乏。
        \item 预防：食物制作应细软，保证荤素搭配，必要时制作成泥状。
    \end{itemize}
\end{enumerate}
\end{defbox}

% =====================================================================
\section{运动员营养}
% =====================================================================

\subsection{运动员营养代谢特点}

\begin{defbox}
\textbf{运动员五大系统生理改变：}

\begin{enumerate}[leftmargin=*]
    \item \textbf{心血管系统：}血容量增加，心输出量大幅提升，供氧供能需求暴涨。
    \item \textbf{神经系统：}超负荷训练易导致神经调节紊乱、疲劳、运动能力下降。
    \item \textbf{消化系统：}运动时胃肠道供血减少，营养素吸收减弱。
    \item \textbf{免疫系统：}长期高强度训练 → 免疫力下降、易呼吸道感染；中低强度有氧运动可提升免疫功能。
    \item \textbf{内分泌系统：}长期大强度训练可致激素紊乱；女性运动员多见月经不调、闭经。
\end{enumerate}

\textbf{整体代谢核心特点：}
\begin{itemize}[leftmargin=*]
    \item 能量消耗极高，易出现氧债，恢复期消耗仍大。
    \item 蛋白质分解增强、尿氮流失多，易发生运动性贫血和中枢疲劳。
    \item 高强度运动时脂肪氧化受限；糖原消耗快，是耐力短板。
    \item 大量出汗，水、钠钾钙镁等矿物质及水溶性维生素流失严重。
    \item 训练刺激微量营养素代谢，需求量高于普通人。
\end{itemize}
\end{defbox}

\subsection{运动员各类营养素需要}

\begin{defbox}
\textbf{1. 能量：}
\begin{itemize}[leftmargin=*]
    \item 总能量：3500-4400 kcal/d（210-280 kJ/kg）。
    \item 供能配比：蛋白质:脂肪:碳水化合物 = 1:(0.7-0.8):4。
    \item 影响因素：运动项目、强度、时长、体重、膳食结构。
\end{itemize}

\textbf{2. 蛋白质：}
\begin{itemize}[leftmargin=*]
    \item 适宜摄入：1.2-2.0 g/(kg$\cdot$d)，优质蛋白占1/3以上。
    \item 作用：修复肌肉、促进恢复、抗疲劳、预防运动性贫血。
    \item 摄入不足 → 耐力下降、乏力；过量加重肝肾负担。
\end{itemize}

\textbf{3. 脂肪：}
\begin{itemize}[leftmargin=*]
    \item 供能占总能量25\%-30\%；耐力项目可至50\%-60\%，极限可达70\%。
    \item 限制饱和脂肪，防血脂升高、降低耐力。
\end{itemize}

\textbf{4. 碳水化合物（核心供能物质）：}
\begin{itemize}[leftmargin=*]
    \item 高强度短时运动：碳水为主；长时间低强度才动用脂肪和蛋白质。
    \item \textbf{糖原填充法：}赛前1周低碳水耗糖原，赛前3天高碳水储备糖原（占膳食70\%）。
    \item 作用：快速供能、无代谢废物，决定耐力上限。
\end{itemize}

\textbf{5. 水与矿物质：}
\begin{itemize}[leftmargin=*]
    \item \textbf{补水：}运动大量出汗，少量多次补水，忌一次性暴饮；补液总量应大于出汗量，同步补电解质。
    \item 钠：高温训练流失多，低钠可致乏力恶心，食盐$<$5 g/d。
    \item 钾：流失致心律失常，推荐3-4 g/d，蔬果补充。
    \item 镁：多汗缺镁 → 抽筋，推荐400-500 mg/d。
    \item 钙：高强度+高温易缺钙、骨质疏松，推荐1000-1500 mg/d。
    \item 铁：训练流失多，易发运动性贫血；女性重点补铁，动物红肉补铁。
    \item 锌：大训练量锌流失、免疫力下降，推荐25 mg/d。
    \item 硒：提升免疫，推荐50-150 $\mu$g/d。
    \item 碘：150 $\mu$g/d，海产品补充。
\end{itemize}

\textbf{6. 维生素（需求量高于普通人群）：}
\begin{itemize}[leftmargin=*]
    \item 维生素B$_{1}$：3-5 mg/d，缺则丙酮酸堆积、神经疲劳；高碳水训练需增量。
    \item 维生素B$_{2}$：2-2.5 mg/d，缺乏损害有氧/无氧运动能力。
    \item 烟酸、B$_{6}$、B$_{12}$、叶酸：参与能量代谢、造血，缺乏致耐力下降。补充时机：赛前2-3周补制剂。
    \item 维生素C：训练推荐200 mg/d；促胶原修复、促铁吸收、抗氧化。
    \item 维生素A：视力需求高项目（射击、击剑）1800 $\mu$g RAE/d。
    \item 维生素D：10-12.5 $\mu$g/d；保护骨骼，缺乏易运动骨折。
    \item 维生素E：高原/耐力项目15-50 mg $\alpha$-TE/d；抗氧化、减少乳酸堆积；不可过量（蓄积中毒）。
\end{itemize}
\end{defbox}

% =====================================================================
\section{特殊环境人员营养}
% =====================================================================

\begin{tipbox}
\textbf{高温环境人员营养需要和膳食原则（了解）：}

\begin{enumerate}[leftmargin=*]
    \item \textbf{水和矿物质：}
    \begin{itemize}[leftmargin=*]
        \item 高温环境大量出汗 → 水和电解质大量丢失（钠、钾、钙、镁）。
        \item 少量多次补水，忌一次性暴饮；同时补充电解质（含钠、钾的运动饮料或淡盐水）。
        \item 食盐摄入比常温环境适当增加（通过膳食或饮水中补充）。
    \end{itemize}

    \item \textbf{能量：}高温环境基础代谢率升高（体温每升高1$^\circ$C，BMR增加约13\%），能量需要较常温增加。
    \item \textbf{蛋白质：}高温下蛋白质分解代谢增强，尿氮排出增加，蛋白质需要量相应提高（优质蛋白占1/2以上）。
    \item \textbf{维生素：}
    \begin{itemize}[leftmargin=*]
        \item 水溶性维生素（B族、C）随汗液流失显著增加。
        \item 维生素C推荐摄入提高至150-200 mg/d。
        \item 维生素B$_{1}$、B$_{2}$按能量比例相应增加（每1000 kcal：B$_{1}$ 0.5-0.6 mg，B$_{2}$ 0.55-0.65 mg）。
    \end{itemize}
    \item \textbf{膳食原则：}清淡可口，增进食欲；多提供汤类和液体食物；保证蔬菜、水果供给充足（补充水和矿物质）；适当增加食盐摄入。
\end{enumerate}

\textbf{低温环境人员营养需要和膳食原则（了解）：}

\begin{enumerate}[leftmargin=*]
    \item \textbf{能量：}低温环境下基础代谢率升高，机体需要额外产热以维持体温 → 能量需要显著增加（可比常温增加10\%-15\%以上）。
    \item \textbf{宏量营养素：}
    \begin{itemize}[leftmargin=*]
        \item 脂肪供能比可提高至35\%-40\%（增加能量密度和饱腹感）。
        \item 碳水化合物供能比维持在50\%-55\%（保证快速供能）。
        \item 蛋白质供给充足（优质蛋白占1/2以上），防止负氮平衡。
    \end{itemize}
    \item \textbf{维生素：}
    \begin{itemize}[leftmargin=*]
        \item 维生素C有助增强耐寒能力，推荐摄入量适当增加。
        \item 维生素B族需要量随能量增加而相应增加。
        \item 维生素A、D、E的摄入也应适当增加。
    \end{itemize}
    \item \textbf{矿物质：}寒冷气候下尿钠排出增加，钠需要量可适当提高；钙、铁、锌、碘的供给亦需充足。
    \item \textbf{膳食原则：}提供高能量密度食物，保证热食热饮供应；增加脂肪摄入比例；食物温热可口；保证蔬菜、水果；食盐摄入可适当增加。
\end{enumerate}
\end{tipbox}

% =====================================================================
\section{复习题}
% =====================================================================

\reviewcontent{
\section{复习题}

\subsection*{一、生命早期1000天与DOHaD}

\begin{enumerate}[leftmargin=*, itemsep=8pt]
    \item \textbf{【名词解释】}生命早期1000天
    \item \textbf{【名词解释】}DOHaD理论
    \item \textbf{【简答题】}简述DOHaD理论的核心含义及其"节俭表型"假说。
\end{enumerate}

\subsection*{二、孕妇营养}

\begin{enumerate}[leftmargin=*, itemsep=8pt, start=4]
    \item \textbf{【名词解释】}生理性贫血
    \item \textbf{【简答题】}列出孕期各阶段的能量和蛋白质增加量。
    \item \textbf{【简答题】}简述叶酸补充对孕期的重要性——为什么应从孕前3个月开始补充？
    \item \textbf{【简答题】}根据孕前BMI分类，列出各组孕期推荐体重增长范围及中晚期每周增长速率。
    \item \textbf{【简答题】}孕期碳水化合物最低摄入量为何不应低于130 g/d？
    \item \textbf{【简答题】}列出《备孕和孕期妇女膳食指南（2022年）》6条关键推荐。
    \item \textbf{【简答题】}简述孕期的内分泌变化——HCG、HCS、雌激素、孕激素各有什么作用和出现时间？
    \item \textbf{【简答题】}列举孕期常见营养问题及预防（至少4项）。
\end{enumerate}

\subsection*{三、乳母营养}

\begin{enumerate}[leftmargin=*, itemsep=8pt, start=12]
    \item \textbf{【简答题】}简述母乳三阶段（初乳、过渡乳、成熟乳）的主要成分差异。
    \item \textbf{【简答题】}列举母乳喂养的八大优点。
    \item \textbf{【简答题】}乳母的能量和蛋白质较非孕期分别增加多少？钙的需求有何特殊性？
    \item \textbf{【简答题】}列举乳母常见营养问题及预防（至少3项）。
\end{enumerate}

\subsection*{四、婴幼儿营养}

\begin{enumerate}[leftmargin=*, itemsep=8pt, start=16]
    \item \textbf{【名词解释】}回应式喂养、辅食
    \item \textbf{【简答题】}简述0-6月龄婴儿母乳喂养指南（2022）的6条准则。
    \item \textbf{【简答题】}简述7-24月龄婴幼儿喂养指南（2022）的6条准则。
    \item \textbf{【简答题】}简述辅食添加的意义（至少4点）及添加原则。
    \item \textbf{【简答题】}简述婴幼儿的铁营养特点——为何6m-2y易患缺铁性贫血？
    \item \textbf{【简答题】}母乳中乳糖的作用是什么？
    \item \textbf{【简答题】}列举婴幼儿常见营养问题及预防（至少5项）。
\end{enumerate}

\subsection*{五、学龄前儿童营养}

\begin{enumerate}[leftmargin=*, itemsep=8pt, start=23]
    \item \textbf{【简答题】}简述学龄前儿童的主要生理特点。
    \item \textbf{【简答题】}列出学龄前儿童膳食指南（2022）核心推荐。
    \item \textbf{【简答题】}学龄前儿童每日奶类、水和体育活动的推荐量分别是多少？
    \item \textbf{【简答题】}列举学龄前儿童常见营养问题及预防。
\end{enumerate}

\subsection*{六、学龄儿童与青少年营养}

\begin{enumerate}[leftmargin=*, itemsep=8pt, start=27]
    \item \textbf{【简答题】}简述学龄儿童骨骼发育的特点。
    \item \textbf{【简答题】}列出学龄儿童膳食指南（2022）5条核心推荐。
    \item \textbf{【简答题】}简述生长发育的4个规律。
    \item \textbf{【简答题】}列举学龄儿童常见营养问题及预防（至少3项）。
    \item \textbf{【简答题】}简述青少年体格发育的特点（身高、体重、形态）。
    \item \textbf{【简答题】}青春期为何是峰值骨量积累的关键期？钙摄入不足的远期后果是什么？
    \item \textbf{【简答题】}为什么青春期女性是缺铁性贫血高发人群？
    \item \textbf{【简答题】}列举青少年常见营养问题及预防（至少4项）。
\end{enumerate}

\subsection*{七、老年人营养}

\begin{enumerate}[leftmargin=*, itemsep=8pt, start=35]
    \item \textbf{【名词解释】}少肌症（Sarcopenia）、合成抵抗（Anabolic Resistance）
    \item \textbf{【简答题】}老年人蛋白质需求量"不减反增"的原因是什么？推荐摄入量是多少？
    \item \textbf{【简答题】}列出《中国老年人膳食指南》（2022）4条条目。
    \item \textbf{【简答题】}列举老年人常见营养问题及预防（至少5项）。
    \item \textbf{【论述题】}从"生命早期1000天 → 成年期 → 老年期"的时间维度，论述生命周期不同阶段营养与健康的连续联系，特别是DOHaD理论和老年人少肌症的关键节点。
\end{enumerate}

\subsection*{八、运动员营养与特殊环境人员营养}

\begin{enumerate}[leftmargin=*, itemsep=8pt, start=40]
    \item \textbf{【简答题】}简述运动员的营养代谢特点（五大系统生理改变及整体代谢特点）。
    \item \textbf{【简答题】}简述运动员各类营养素需要要点（碳水化合物、蛋白质、水与矿物质）。
    \item \textbf{【简答题】}简述高温和低温环境人员的营养需要和膳食原则。
\end{enumerate}

\subsection*{参考答案要点}

\begin{enumerate}[leftmargin=*, itemsep=5pt]
    \item \textbf{生命早期1000天：}孕期270天 + 0-6月龄180天（母乳喂养）+ 7-24月龄（辅食添加）。是生长发育的关键期和敏感期，营养状态对终身健康产生深远和不可逆的影响。
    \item \textbf{DOHaD理论（健康与疾病的发育起源）：}生命早期（胎儿、婴幼儿期）经历不利因素（如营养不良、环境不良等），将会增加其成年后患肥胖、糖尿病、心血管疾病等慢性疾病的几率。核心机制为发育程序化（表观遗传调控）。
    \item DOHaD核心含义：早期营养环境通过表观遗传机制（DNA甲基化、组蛋白修饰）对远期疾病风险产生持久性影响。节俭表型假说：宫内营养不良使胎儿产生"预测性适应"，若出生后实际面临"营养过剩"环境，则代谢-环境不匹配大幅增加成年期代谢综合征风险。
    \item \textbf{生理性贫血：}孕期血容量增加，血浆容积增加幅度 $>$ 红细胞增加幅度，导致血液稀释。通常在孕20-30周最明显。WHO标准：孕期Hb $\geq$ 110 g/L，孕中期 $\geq$ 105 g/L。
    \item 孕早期：E+0，P+0；孕中期：E+250 kcal/d，P+15 g/d；孕晚期：E+400 kcal/d，P+30 g/d。
    \item 叶酸缺乏致神经管缺陷（NTDs）。孕后3-8周是神经管闭合关键期（此时许多妇女尚不知已怀孕），故需孕前3个月开始补充。推荐量：400 $\mu$g DFE/d。
    \item 低体重（BMI$<$18.5）：总11.0-16.0 kg，每周0.46 kg；正常（18.5-24.0）：总8.0-14.0 kg，每周0.37 kg；超重（24.0-28.0）：总7.0-11.0 kg，每周0.30 kg；肥胖（$\geq$28.0）：总5.0-9.0 kg，每周0.22 kg。
    \item 防止酮体对胎儿中枢神经系统的毒性。严重早孕反应不能进食时需静脉补充葡萄糖。
    \item 6条关键推荐：(1)调整孕前体重至正常范围，保证孕期体重适宜增长；(2)常吃含铁丰富的食物，选用碘盐，合理补充叶酸和维生素D；(3)孕吐严重者，可少量多餐，保证摄入含必需量碳水化合物的食物；(4)孕中晚期适量增加奶、鱼、禽、蛋、瘦肉的摄入；(5)经常户外活动，禁烟酒，保持健康生活方式；(6)愉快孕育新生命，积极准备母乳喂养。
    \item HCG：受精1周开始出现，8-9周达高峰，维持妊娠早期孕酮分泌。HCS：妊娠6周出现，36周达高峰至分娩，促进乳腺发育和胎儿生长，降低胰岛素敏感性。雌激素：分泌增加，增加子宫和胎盘血流量，促进乳房发育。孕激素：分泌增加，减少子宫收缩，维持妊娠。
    \item 孕期常见营养问题：(1)早孕反应——HCG升高+贲门松弛，少量多餐保证碳水130g/d；(2)孕期便秘——孕酮致肠蠕动减慢+子宫压迫；(3)GDM——生理性胰岛素抵抗，控制体重和低GI饮食；(4)妊娠高血压——与钙不足可能有关，保证钙800mg/d；(5)缺铁性贫血——血容量扩张+胎铁转运，孕中25mg/d、孕晚29mg/d铁；(6)钙营养——钙吸收率代偿升高，800mg/d即可。
    \item 初乳（第1周）：淡黄色，蛋白质多，富含免疫球蛋白（sIgA和乳铁蛋白），乳糖和脂肪相对较少。过渡乳（第2周）：蛋白质$\downarrow$，乳糖和脂肪$\uparrow$。成熟乳（第2周以后）：乳白色，蛋白质相对较少，乳糖和脂肪多。
    \item 八大优点：(1)乳清蛋白为主，必需氨基酸比例适当；(2)脂肪颗粒小+乳脂酶→易消化吸收；(3)富含乳糖→促进乳酸杆菌、抑制大肠埃希菌，利于钙铁锌吸收；(4)钙磷比适宜（2:1）；(5)含大量免疫物质；(6)不易过敏；(7)经济方便卫生；(8)促进产后恢复，增进母婴交流。
    \item 能量：+400 kcal/d。蛋白质：+25 g/d。钙：RNI 800 mg/d（不增加），乳汁中钙含量较稳定（每日排出约300 mg），乳母钙供给不足时会动用自身骨骼中的钙满足乳汁钙含量——哺乳期骨骼密度下降，断奶后可逆性恢复。
    \item 乳母常见营养问题3项：(1)能量和营养素摄入不足——产后过度节食+饮食不规律→泌乳量减少，能量+400kcal/d蛋白+25g/d；(2)骨骼健康——乳汁每日排钙约300mg，钙不足动用骨钙，保证800mg/d；(3)产后贫血——分娩失血+哺乳消耗，铁不能通过乳腺入乳，母亲需注意补铁。
    \item \textbf{回应式喂养：}按婴儿饥饿信号适时喂养，而非固定时间刻板喂养。\textbf{辅食：}除母乳外，为满足婴幼儿生长发育需要而添加的其他食物。
    \item 6条准则：(1)母乳是婴儿最理想的食物，坚持6月龄内纯母乳喂养；(2)生后1小时内开奶；(3)回应式喂养；(4)适当补充维生素D，母乳喂养无需补钙；(5)任何动摇母乳喂养的想法和举动都必须咨询医生；(6)定期监测体格指标。
    \item 6条准则：(1)继续母乳喂养，满6月龄起必须添加辅食，从富含铁的泥糊状食物开始；(2)及时引入多样化食物，重视动物性食物的添加；(3)尽量少加糖、盐，油脂适当，保持食物原味；(4)提倡回应式喂养，鼓励但不强迫进食；(5)注重饮食卫生和进食安全；(6)定期监测体格指标，追求健康生长。
    \item 辅食添加意义：(1)补充营养素——6月后母乳不能满足全部需要，99\%铁需从辅食获得；(2)促进消化系统发育；(3)培养咀嚼吞咽能力；(4)促进味觉发育和饮食习惯形成；(5)促进神经心理发育。添加原则：每次一种新食物观察3-5天；由少到多由稀到稠由细到粗；健康消化功能正常时添加；不加盐糖；1岁后淡口味。
    \item 足月新生儿有足够铁贮存供4-6个月需要。6月龄时体内储存铁已耗尽，7-12月龄99\%的铁必须从辅食中获得。6m-2y为缺铁性贫血高发期。7-12月龄铁RNI = 10 mg/d。
    \item 乳糖促进乳酸杆菌生长、抑制大肠埃希菌生长；利于钙、铁、锌吸收。
    \item 婴幼儿常见营养问题5项：(1)缺铁性贫血——6m-2y高发，补充富铁辅食；(2)VD缺乏性佝偻病——生后数开始补VD 400IU/d；(3)PEM——辅食不及时质差；(4)锌缺乏——生长迟缓食欲差；(5)碘缺乏——克汀病，使用碘盐；(6)VK缺乏——新生儿肌注VK；(7)龋齿——1岁前不加糖不喝含糖饮料。
    \item (1)每年体重增长约2 kg，身高约5-7 cm；(2)神经系统仍属脑细胞生长发育关键期；(3)3岁乳牙已出齐，6岁恒牙已萌出，但咀嚼和消化能力远低于成人；(4)注意力分散（约15分钟），无法专心进食。
    \item (1)食物多样，规律就餐，自主进食；(2)每天饮奶（300-500 mL），足量饮水（600-800 mL），合理选择零食；(3)合理烹调，少调料少油炸；(4)参与食物选择制作；(5)经常户外活动，定期体格测量。
    \item 奶类每日300-500 mL；水每日600-800 mL；每天至少60分钟体育活动；每天视屏时间不超过2小时。
    \item 学龄前儿童常见营养问题：(1)挑食偏食——家长以身作则多次尝试；(2)零食过多——选健康零食两餐之间；(3)含糖饮料——白开水为主；(4)钙不足——奶类300-500ml/d；(5)铁锌不足——保证肉类肝脏；(6)超重肥胖——控能量+户外活动60min/d。
    \item 骨骼特点：骨胶多、钙盐少弹性大、硬度小，不易骨折而易弯曲变形。
    \item (1)主动参与食物制作和选择，提高营养素养；(2)吃好早餐，合理选择零食，培养健康饮食行为；(3)天天喝奶，足量饮水，不喝含糖饮料，禁止饮酒；(4)多户外活动，少视屏时间，每天60分钟以上中高强度身体活动；(5)定期监测体格发育，保持体重适宜增长。
    \item 生长发育4规律：(1)连续性；(2)不同步；(3)波浪式；(4)差异性。
    \item 学龄儿童常见营养问题：(1)不吃早餐或早餐质量差——早餐三类食物以上；(2)零食甜饮料过度——以白开水为日常饮水；(3)钙和VD不足——天天喝奶+户外活动；(4)超重肥胖或消瘦两极分化——平衡膳食+每天60min中高强度活动；(5)脊柱骨骼问题——保证钙摄入+正确坐姿+运动。
    \item 身高：性别差异——男较女晚2年，男7-9 cm/y，女5-7 cm/y。体重：4-5 kg/y，时间较长。形态：四肢比脊柱早，脚最先，从远到近。
    \item 青春期是获得峰值骨量的最关键阶段，约45\%的成人骨量在青春期积累。钙摄入不足 → 无法达到遗传决定的峰值骨量 → 老年期骨质疏松风险大幅增加。钙RNI：1000-1200 mg/d。
    \item 青春期女性需弥补月经铁的丢失，同时生长发育需要大量铁。11-13岁：18 mg/d；14-17岁：18 mg/d。是缺铁性贫血高发人群。
    \item 青少年常见营养问题：(1)青春期女性IDA——月经铁丢失+生长，铁RNI 18mg/d；(2)钙不足影响峰值骨量——1000-1200mg/d；(3)不健康减肥——盲目节食致能量和营养素不足影响发育；(4)不吃早餐和在外就餐——油盐多蔬果少；(5)超重肥胖——控能量+60min中高强度运动；(6)锌碘缺乏——锌9-11.5mg/d碘110-120$\mu$g/d；(7)情绪化进食/饮食行为偏离。
    \item \textbf{少肌症：}50岁后骨骼肌量平均每年减少1\%-2\%，60岁以上慢性肌肉丢失约30\%，80岁以上为50\%。表现为骨骼肌质量、力量和功能的进行性丧失。\textbf{合成抵抗：}衰老肌肉对较低剂量氨基酸/蛋白质的合成代谢反应下降，需要更多蛋白质刺激同等幅度的肌蛋白合成。
    \item 原因：合成抵抗（Anabolic Resistance）——老年人蛋白质合成对摄入蛋白质的敏感性下降，分解代谢亢进。推荐摄入量：1.0-1.5 g/(kg$\cdot$bw)/d。65岁以上：男72 g/d，女62 g/d（高于50岁组的65和55 g/d）。
    \item (1)食物品种丰富，动物性食物充足，常吃大豆制品；(2)鼓励共同进餐，保持良好食欲；(3)积极户外活动，延缓肌肉衰减，保持适宜体重；(4)定期健康体检，测评营养状况，预防营养缺乏。
    \item 列举老年人常见营养问题的5项以上：(1)少肌症——50岁后骨骼肌年均减1\%-2\%，蛋白质1.0-1.5 g/(kg$\cdot$d)+抗阻运动；(2)PEM——食欲降低、咀嚼困难致营养不足，保证动物性食物和共同进餐；(3)骨质疏松——钙+VD+负重运动；(4)VD缺乏——晒太阳+65岁以上RNI 15 $\mu$g/d；(5)便秘——增加膳食纤维和饮水；(6)隐性脱水——渴觉减退，每日饮水不少于1500 ml；(7)VB$_{12}$缺乏——萎缩性胃炎致吸收障碍；(8)牙齿口腔问题——食物细软。
    \item \textbf{论述题答题要点：}
    \begin{itemize}[leftmargin=*]
        \item \textbf{生命早期1000天 → DOHaD：}孕期270d + 0-6月龄180d + 7-24月龄为机遇窗口。早期营养不良经表观遗传（发育程序化）导致远期代谢综合征高风险。节俭表型假说——宫内"预测匮乏"+出生后"营养过剩"→ 代谢失匹配。
        \item \textbf{成年期：}膳食行为和生活方式继续强化或修正早期程序化效应。膳食指南2022八准则贯穿全生命周期。
        \item \textbf{老年期 → 少肌症与合成抵抗：}蛋白质需要量不减反增（1.0-1.5 g/(kg$\cdot$bw)/d），能量需求下降但蛋白质需求上升形成"老年人营养悖论"。需配合抗阻运动以实现成功衰老。
        \item \textbf{核心结论：}营养影响不是点状而是连续的——早期干预效费比最高。从生命早期1000天到老年期，营养贯彻生命周期理念。
    \end{itemize}

    \item 运动员五大系统生理改变：心血管（血容量和心输出量增加）、神经（超负荷易疲劳紊乱）、消化（运动时供血减少影响吸收）、免疫（高强度训练致免疫下降，中低强度有氧提升免疫）、内分泌（大强度可致激素紊乱，女性月经不调闭经）。整体代谢特点：能量消耗极高，蛋白分解增强尿氮流失多，糖原消耗快为耐力短板，大量出汗致水和矿物质维生素流失，微量营养素需求高于常人。
    \item 运动员营养素需要要点——碳水化合物：高强度短时运动核心供能，糖原填充法（赛前1周低碳水耗糖原，赛前3天高碳水储备占膳食70\%）；蛋白质：1.2-2.0 g/(kg$\cdot$d)，优质蛋白1/3以上；水与矿物质：少量多次补液，总补液量大于出汗量同步补电解质；钠钾钙镁铁锌硒碘均需充分补充。
    \item 高温环境人员：大量出汗致水电解质流失，少量多次补水+补电解质；BMR升高能量增加；蛋白分解增强；水溶性维生素随汗液流失显著增加（VC 150-200 mg/d）；清淡可口膳食多汤类多蔬果。低温环境人员：BMR升高能量显著增加（10\%-15\%以上）；脂肪供能比可提高至35\%-40\%；碳水50\%-55\%；蛋白充足；VC、B族增加；食盐可适当增加；保证热食热饮，食物温热可口。
\end{enumerate}

\newpage
}
